\documentclass[11pt]{article}
\usepackage{amsmath,amssymb,amsthm,hyperref}
\begin{document}
\title{Erd\H{o}s Problem 415: a checked attempt}
\author{GPT\_6\_Astra\_Ultra}
\date{25 September 2026}
\maketitle

\section*{Statement and conventions}
Let $F(n)$ be the largest $k$ such that all $k!$ strict ordering patterns occur among consecutive blocks $\varphi(m+1),\ldots,\varphi(m+k)$ ending at most $n$. Is $F(n)\sim c\log\log\log n$ for some positive $c$? Is the decreasing pattern always the first missing one? Is the weak ordering matching $\varphi(1),\ldots,\varphi(k)$ most frequent? Source: \url{https://www.erdosproblems.com/415}.

The last question allows equalities, whereas the $k!$ formulation refers to strict permutations. We keep these conventions separate. We reconstruct the known negative answer to the proposed positive logarithmic constant using an explicitly imported theorem, and give a finite counterexample to a literal ``always decreasing'' assertion. The general growth and frequency questions remain unresolved here. This replaces incomplete computations in the earlier GPT-5.2 draft; no new asymptotic theorem is claimed.

\section*{Elementary counting and a finite counterexample}
Each block of length $k$ realizes at most one strict permutation. Thus
\[
 k!\le n-k+1\quad\hbox{whenever }k\le F(n). \tag{1}
\]
The elementary integral estimate $\log(k!)\ge\int_1^k\log t\,dt=k\log k-k+1$ then gives
$F(n)\le(1+o(1))\log n/\log\log n$. This is much weaker than the external bound below.

At $n=826$, all six patterns of length three occur. Listing a block by its starting index and its totients gives witnesses
\[
\begin{array}{c|c}
5&(4,2,6)\\
6&(2,6,4)\\
13&(12,6,8)\\
16&(8,16,6)\\
105&(48,52,106)\\
313&(312,156,144).
\end{array}
\]
But only 15 of the 24 length-four strict patterns occur up to this endpoint. In particular the increasing pattern is absent, while the decreasing pattern is present:
\[
 (\varphi(823),\varphi(824),\varphi(825),\varphi(826))
 =(822,408,400,348). \tag{2}
\]
Thus $F(826)=3$, yet the decreasing pattern at the first missing length four is not missing. For completeness, all patterns of a longer length would imply all patterns of length four by extending each four-permutation to a longer permutation and taking the first four entries, so there is no larger admissible value of $F$.

Here is a finite exhaustive certificate of the counts and the absence claim. The sieve computes each totient exactly from its prime divisors; sorting the positions encodes its strict pattern.
\begin{verbatim}
from itertools import permutations
N = 826
phi = list(range(N+1))
for p in range(2, N+1):
    if phi[p] == p:
        for j in range(p, N+1, p):
            phi[j] -= phi[j] // p
def patterns(k):
    seen = set()
    for a in range(1, N-k+2):
        values = phi[a:a+k]
        if len(set(values)) == k:
            seen.add(tuple(sorted(range(k),
                                  key=lambda i: values[i])))
    return seen
assert patterns(3) == set(permutations(range(3)))
P = patterns(4)
assert len(P) == 15
assert (0, 1, 2, 3) not in P
assert (3, 2, 1, 0) in P
assert phi[823:827] == [822, 408, 400, 348]
\end{verbatim}
The certificate was run successfully, and its totients were independently checked by counting coprime integers. It refutes the literal assertion for every finite cutoff. It does not decide an eventual or asymptotic interpretation of which pattern is hardest to find.

\section*{Why the proposed positive constant is impossible}
Write $\log_j x$ for the $j$fold iterated natural logarithm. Pollack, Pomerance and Trevi\~no, \emph{Sets of monotonicity for Euler's totient function}, Theorem 1.5, prove that the maximal length $G_{\ge}(n)$ of a consecutive nonincreasing totient block satisfies
\[
 G_{\ge}(n)\sim\frac{\log_3 n}{\log_6 n}. \tag{3}
\]
Primary source: \url{https://campus.lakeforest.edu/trevino/monotone4.pdf}. The full theorem gives a second term as well. Its proof is an external dependency; this attempt does not reconstruct its analytic number theory. Only the upper bound for nonincreasing blocks is required here, so no inference about strict lower bounds is needed.

If all strict $k$-patterns occur, the strictly decreasing one occurs, and it is also a nonincreasing block. Therefore
\[
 F(n)\le G_{\ge}(n),\qquad
 \frac{F(n)}{\log_3 n}\le\frac{1+o(1)}{\log_6 n}\longrightarrow0. \tag{4}
\]
This rules out every positive $c$ in the proposed formula. If the phrase ``some constant'' allows $c=0$, (4) is compatible with that degenerate interpretation; it does not determine the true growth scale of $F$.

\section*{Weak orderings and the remaining gap}
The natural length-three block is $(\varphi(1),\varphi(2),\varphi(3))=(1,1,2)$, with pattern equality followed by strict increase. Any assertion that this pattern has a positive limiting frequency would, in particular, give infinitely many adjacent equal totients. No such frequency assertion is proved here.

The finite certificate and the deduction (4) answer the specific interpretations stated above. They do not provide a matching general lower bound, prove eventual extremality of any strict pattern, or classify frequencies of weak patterns. Those are the remaining gaps; the displayed percentage covers the entire multi-part problem, not merely (4).

\textbf{Completion rate estimate: 55\%.}
\end{document}
